\documentclass[12pt, a4paper]{article}

\usepackage[UTF8]{ctex}              % 中文支持，关键
\usepackage[margin=1.5cm]{geometry}  % 页边距减小到原来约一半
\usepackage{mdframed}
\usepackage{stmaryrd}
\usepackage{amssymb}
\usepackage{amsmath}

\begin{document}

\title{离散数学笔记}
\author{DuRuize}
\date{}
\maketitle
\setlength{\parindent}{0pt}

\newpage
\section*{2026-9-10}
\subsection*{0.先来分清两个概念}
\textbf{最令人疑惑的点是:} $\models$ 和 $\vdash$ 的区别.

\textbf{语法后承：} $\Gamma \vdash \varphi$ 是指 $\dfrac{\psi \qquad \psi\to\varphi}{\varphi}$ 这样的能推理出的过程

\textbf{语义后承：} $\Gamma \models \varphi$ 是指 $\Bigl(\forall \psi \in \Gamma,\ \llbracket \psi \rrbracket_v = 1\Bigr) \Rightarrow \llbracket \varphi \rrbracket_v = 1$

\subsection*{1. 完备性定理:}

\textbf{语义后承：} $\Gamma \models \phi$

\textbf{soundness（若可证明则可语义推出）：} $\Gamma \vdash \phi \Rightarrow \Gamma \models \phi$

\textbf{completeness（若可语义推出则可证明）：} $\Gamma \models \phi \Rightarrow \Gamma \vdash \phi$

\textbf{我们接下来引入若干概念就是为了证明上面两点.}

\subsection*{2. 一致集与扩充成极大一致集}
\vspace{0.3cm}
\textbf{一致性(consistent)：}
\begin{itemize}
    \item 若 $\Gamma \nvdash \bot$，则称 $\Gamma$ 是一致的
    \item 若 $\Gamma \vdash \bot$，则称 $\Gamma$ 是不一致的
\end{itemize}

\textbf{极大一致性（maximally consistent）：}
\begin{itemize}
    \item $\Gamma$ 是一致的
    \item $\forall \, \Gamma' \supset \Gamma, \Gamma'$ 是不一致的
\end{itemize}
\textit{性质：} 若 $\Gamma$ 极大一致，则 $\forall \phi, \Gamma \vdash \phi \text{ 或 } \Gamma \vdash \neg \phi$

\vspace{0.3cm}
\textbf{演绎封闭（close under derviation）：}
$\forall \phi, \text{ 若 } \Gamma \vdash \phi \text{，则 } \phi \in \Gamma$


\textbf{引理：} 对任意一致集 $\Gamma$，存在极大一致集 $\Gamma'$，使得 $\Gamma \subseteq \Gamma'$。

\textbf{证明：}
\begin{enumerate}
    \item 由于公式集是可数的(每一位都是有限种组合)，枚举所有公式：$\phi_0, \phi_1, \phi_2, \dots$
    \item 递归构造集合序列 $\{\Gamma_n\}_{n \in \mathbb{N}}$：
    \begin{itemize}
        \item $\Gamma_0 = \Gamma$
        \item $\Gamma_{n+1} = 
        \begin{cases} 
        \Gamma_n \cup \{\phi_n\}, & \text{若 } \Gamma_n \cup \{\phi_n\} \text{ 一致} \\
        \Gamma_n \cup \{\neg \phi_n\}, & \text{否则}
        \end{cases}$
    \end{itemize}
    \item 定义 $\Gamma' = \bigcup\limits_{n=0}^{\infty} \Gamma_n$ \text{则} $\Gamma'$ 一致且极大。
\end{enumerate}

\subsection*{3.证明$\Gamma \text{ 一致} \Rightarrow \exists v \ \  s.t. \forall \psi \in \Gamma, \llbracket \psi \rrbracket_v = 1$。}

\textbf{证明：}
\begin{enumerate}
    \item 由上述引理，将 $\Gamma$ 扩充为极大一致集 $\Gamma'$。
    \item 基于 $\Gamma'$ 定义赋值 $v$：
    \[
    \llbracket P_i \rrbracket_v = 
    \begin{cases} 
    1, & \text{若 } P_i \in \Gamma' \\
    0, & \text{若 } \neg P_i \in \Gamma'
    \end{cases}
    \]
    \item 对公式 $\psi$ 进行\textbf{结构归纳}，证明：
    \[
    \llbracket \psi \rrbracket_v = 1 \iff \psi \in \Gamma'
    \]
    \begin{itemize}
        \item \textit{基础情形：} 原子命题 $P_i$，由 $v$ 的定义成立。
        \item \textit{归纳步骤：}
        \begin{itemize}
            \item 情形 $\psi = \neg \phi$：
            $\llbracket \neg \phi \rrbracket_v = 1 \iff \llbracket \phi \rrbracket_v = 0 \iff \phi \notin \Gamma' \iff \neg \phi \in \Gamma'$（由极大性）。
            \item 情形 $\psi = \phi_1 \to \phi_2$：
            $\llbracket \phi_1 \to \phi_2 \rrbracket_v = 0 \iff \llbracket \phi_1 \rrbracket_v = 1 \text{ 且 } \llbracket \phi_2 \rrbracket_v = 0$。
            由归纳假设，这意味着 $\phi_1 \in \Gamma'$ 且 $\phi_2 \notin \Gamma'$。
            由于 $\Gamma'$ 演绎封闭，若 $\phi_1 \to \phi_2 \in \Gamma'$，则 $\phi_2 \in \Gamma'$（矛盾）。
            因此 $\phi_1 \to \phi_2 \notin \Gamma'$，所以 $\neg(\phi_1 \to \phi_2) \in \Gamma'$。
            于是 $\llbracket \phi_1 \to \phi_2 \rrbracket_v = 1 \iff \phi_1 \to \phi_2 \in \Gamma'$。
        \end{itemize}
    \end{itemize}
\end{enumerate}

\subsection*{4. 证明 $\vdash \,=\, \models$}

\textbf{目标：} 证明
\[
\Gamma \vdash \varphi \iff \Gamma \models \varphi
\]
即语法推出关系与语义推出关系重合。

\vspace{0.3cm}
\textbf{（1）可靠性：$\Gamma \vdash \varphi \Rightarrow \Gamma \models \varphi$}

对证明长度做归纳：
\begin{itemize}
    \item 若 $\varphi \in \Gamma$，则任何满足 $\Gamma$ 的赋值 $v$ 都满足 $\varphi$。
    \item 若 $\varphi$ 是逻辑公理，则 $\varphi$ 是永真式，任何 $v$ 都满足它。
    \item 若 $\varphi$ 由推理规则得到，例如 MP：
    \[
    \frac{\psi \qquad \psi \to \varphi}{\varphi}
    \]
    由归纳假设，$\psi$ 与 $\psi \to \varphi$ 在任意满足 $\Gamma$ 的 $v$ 下都为真；而真值表给出
    \[
    \llbracket \psi \rrbracket_v = 1,\quad
    \llbracket \psi \to \varphi \rrbracket_v = 1
    \Rightarrow \llbracket \varphi \rrbracket_v = 1.
    \]
    因此 $\varphi$ 也被满足。
\end{itemize}
所以
\[
\Gamma \vdash \varphi \Rightarrow \Gamma \models \varphi.
\]

\vspace{0.3cm}
\textbf{（2）完备性：$\Gamma \models \varphi \Rightarrow \Gamma \vdash \varphi$}

用逆否命题。假设 $\Gamma \nvdash \varphi$，则 $\Gamma \cup \{\neg \varphi\}$ 一致。

于是由第 3 节的结果，存在赋值 $v$，使得
\[
\forall \psi \in \Gamma \cup \{\neg \varphi\},\quad \llbracket \psi \rrbracket_v = 1.
\]
也就是说，对每个 $\psi \in \Gamma$ 都有
\[
\llbracket \psi \rrbracket_v = 1,
\]
并且
\[
\llbracket \neg \varphi \rrbracket_v = 1,
\]
即
\[
\llbracket \varphi \rrbracket_v = 0.
\]
但由假设 $\Gamma \models \varphi$，任何满足 $\Gamma$ 的赋值 $v$ 都必须满足 $\varphi$，即
\[
\Bigl(\forall \psi \in \Gamma,\ \llbracket \psi \rrbracket_v = 1\Bigr)
\Rightarrow
\llbracket \varphi \rrbracket_v = 1.
\]
这与前面得到的 $\llbracket \varphi \rrbracket_v = 0$ 矛盾。

因此假设 $\Gamma \nvdash \varphi$ 不成立，故
\[
\Gamma \vdash \varphi.
\]
所以
\[
\Gamma \models \varphi \Rightarrow \Gamma \vdash \varphi.
\]
\textbf{目前为止是前半节课想要证明的(?)}
\vspace{0.3cm}

\subsection*{5. 经典逻辑与反证法（区分直觉主义与经典逻辑）}

\textbf{反证法(RAA)：} 假设\neg P 推出 \bot
\textbf{无需反证法即可证明的公式（直觉主义逻辑）：}
\begin{itemize}
    \item $P \to \neg\neg P$
    \item $\neg\neg\neg P \to \neg P$
    \item $(P \to Q) \to (\neg Q \to \neg P)$
\end{itemize}

\textbf{需要反证法的公式（经典逻辑）：}
\begin{itemize}
    \item $\neg\neg P \to P$
    \item $(\neg P \to P) \to P$
    \item $((P \to Q) \to P) \to P$ \quad 
\end{itemize}

注意:这里说到的反证法是我们的证明系统内部的RAA,而4.中我们用到的反证法是数学意义上的反证法,不是一个范围内的(这或许类似于不同域上的运算限制?)

\subsection*{6. 一阶逻辑的字母表与项}
\textbf{字母表：}
\begin{itemize}
    \item 变元符号：\(x, y, z, \dots\)
    \item 常元符号：\(c_1, c_2, \dots\)
    \item 函数符号：\(f_1^1, f_1^2, \dots, f_2^1, f_2^2, \dots\) (上标表示元数)
    \item 谓词符号：\(P_1^1, P_1^2, \dots\) 及等词 \(=\)
    \item 联结词：\(\neg, \to, \wedge, \vee, \leftrightarrow\)
    \item 量词：\(\forall, \exists\)
    \item 辅助符号：\((, )\)
\end{itemize}

\textbf{项（Term）：} 归纳定义
\begin{itemize}
    \item 任何变元 \(x\) 是项：\(x \in \text{TERM}\)
    \item 任何常元 \(c\) 是项：\(c \in \text{TERM}\)
    \item 若 \(t_1, \dots, t_k\) 是项，则 \(f_i^k(t_1, \dots, t_k)\) 是项
\end{itemize}
即 \(\text{TERM} \subseteq \text{Alphabet}^*\)。

\subsection*{7. FOR}
\textbf{FOR(formula)：} 归纳定义
\begin{itemize}
    \item 若 \(t_1, \dots, t_k\) 是项，则 \(P_i^k(t_1, \dots, t_k) \in \text{FOR}\)
    \item 若 \(\varphi, \psi \in \text{FOR}\)，则 \(\neg\varphi, \varphi \wedge \psi, \varphi \vee \psi, \varphi \to \psi,\varphi = \psi \in \text{FOR}\)
    \item 若 \(\varphi \in \text{FOR}\) 且 \(x\) 是变元，则 \(\forall x \varphi, \exists x \varphi \in \text{FOR}\)
\end{itemize}

\subsection*{8. 自由变元与约束变元（Free/Bound）}
\textbf{项 \(t\) 的变元集：} \(\text{Var}(t) = \{x \mid x \text{ 是 } t \text{ 中的变元符号}\}\)

\textbf{公式 \(\varphi\) 的变元集：}
\begin{itemize}
    \item \(\varphi = P^k(t_1, \dots, t_k)\)：\(\text{Var}(\varphi) = \text{Free}(\varphi) = \bigcup_{i=1}^k \text{Var}(t_i)\) 注:term不包含量词,所以term的var都是free
    \item \(\varphi = \psi_1 \square \psi_2\)：\(\text{Var}(\varphi) = \text{Var}(\psi_1) \cup \text{Var}(\psi_2)\)
    \item \(\varphi = \forall x \psi\) 或 \(\exists x \psi\)：
    \[
    \text{Var}(\varphi) = \text{Var}(\psi), \quad \text{Free}(\varphi) = \text{Var}(\psi) \setminus \{x\}
    \]
\end{itemize}
\textit{注意：量词所约束的变元称为约束变元（bound），未被约束的称为自由变元（free）。自由变元可以替换成其他独立的记号}

\subsection*{9. 结构（Structure）与模型（Model）}
\textit{结构(structure)：} \(\langle \Omega, R_1, \dots, R_n, f_1, \dots, f_n, c_1, \dots \rangle\)
\begin{itemize}
    \item \(\Omega\) 是非空集合（论域）
    \item 对每个 \(f_i^k\)，有 \(f_i^k: \Omega^k \to \Omega\)
    \item 对每个 \(R_i^k\)，有 \(R_i^k \subseteq \Omega^k\)
\end{itemize}


\subsection*{10. 项和公式的语义解释}
\textbf{项的解释：}
\begin{itemize}
    \item 常元：\([\bar{c_i}]_\mathcal{M} = c_i\)
    \item 复合项：\([f_i^k(t_1, \dots, t_k)]_\mathcal{M} = f_i^k([t_1]_\mathcal{M}, \dots, [t_k]_\mathcal{M})\)
\end{itemize}

\textbf{公式的真值：}
\begin{itemize}
    \item 原子公式：\([P_i^k(t_1, \dots, t_k)]_\mathcal{M} = 1 \iff ([t_1]_\mathcal{M}, \dots, [t_k]_\mathcal{M}) \in P_i^k\)
    \item 量词：\([\forall x \varphi]_\mathcal{M} = \min_{x \in \Omega} [\varphi(x)]_\mathcal{M}\)
    \item \([\exists x \varphi]_\mathcal{M} = \max_{x \in \Omega} [\varphi(x)]_\mathcal{M}\)
\end{itemize}

\subsection*{11. 等词公理}
\begin{itemize}
    \item 自反性：\(\forall x (x = x)\)
    \item 对称性：\(\forall x \forall y (x = y \to y = x)\)
    \item 传递性：\(\forall x \forall y \forall z ((x = y \wedge y = z) \to x = z)\)
\end{itemize}

\end{document}